\section{Review object and substantive changes}
We thank the referee for identifying the difference between learning a better witness and improving an economic policy. R19 responds with a new, executed policy-value mechanism rather than another enclosure of the same actor. The title, original stopped economy, original action set, and full-domain $.01$ objective are unchanged. No inherited actor or optimal-value label enters the new training. Every old mathematical and numerical file remains unchanged, and all substantive R18 sections are retained in the current paper or supplement.

This response concerns the report on the materialized R18 manuscript, not the earlier R16 report. The report is in \path{reviews/2026-09-23-econometrica-r18/referee_report.md}. Its reviewed manuscript commit begins \texttt{49d286a}; the review input commit begins \texttt{074849b}. The new scientific protocol was committed at \texttt{2e19256} before the R19 randomized studies. Full immutable identifiers and file hashes are supplied in the publication manifest.

The principal new evidence is as follows. Five model-only neural time-control generators are trained at four predeclared work levels. Each of the fifteen noninitial proposals passes an independent policy-payoff ordering test before further updates. All final reference-state optimality bounds are below $.01$. A new coupled-payoff theorem certifies that the true initial-to-final policy gain is positive throughout the initial-state rectangle $K=[1.98,2.02]\times[1.24,1.26]$. Combining that gain with a new rectangular regret bound proves a reduction in true regret of at least 58 percent uniformly on $K$, and at least 75 percent at the reference state. Independent evaluation of the same new policies also supplies one-sided continuation-value targets exceeding settlement by more than 6.708 on a preference segment of the inward wealth face.

These are certificates for new neural-generated policies in the unchanged economy. They are not certificates for the old full-state feedback network over the entire time--state cylinder. The old full-domain bound remains approximately 7.278, and we do not call it $.01$. The rectangular bound is $\NineteenK$, also not falsely rounded down to $.01$. The new method therefore provides demonstrated policy learning and accurate reference-state control while keeping the full-domain requirement explicit.

\section{Scientific findings F1--F11}
\subsection{R18-F1 --- Flagship policy accuracy}
\emph{Response.} We introduce a materially new solver object, train it from five fresh random initializations, and independently evaluate every predeclared work level. The actor is a two-hidden-layer, width-16 tanh network with 338 parameters. It generates 16 consumption and preference-control slabs; the portfolio is zero. Consumption is budget normalized, and a checked wealth-transfer rule gives an initial-state-indexed family on $K$. The original admissible control set is not changed: this is a feasible neural-generated subclass evaluated against an upper bound for the unrestricted original problem.

The initial reference-state regret bounds are about $.0392$--$.0398$. Final bounds range from $\NineteenBest$ to $\NineteenWorst$. Every final policy reaches the original numerical threshold at $(t,u,x,k)=(0,2,1.25,2)$. The main table displays all four budgets for every seed. A distinct coupled-payoff proof and four-vertex optimality argument extend policy improvement and a useful regret bound to every state in $K$. We report both generation and independent evaluation costs, including earlier checks at first target attainment.

This addresses the request for new actors, multiple seeds, increasing work, complete evaluation of each retained object, and a policy-sensitive criterion. It does not close the full-cylinder feedback-accuracy gate: the new state-indexed time policies are not relabeled as the old feedback architecture, and the $K$ certificate is not extrapolated to all initial times and states. The full-domain objective remains a substantive target of the paper rather than being removed or replaced by a different economy.

\subsection{R18-F2 --- Witness descent versus policy improvement}
\emph{Response.} We separate the two questions mathematically and experimentally. The new acceptance test requires the new policy's independently verified payoff lower endpoint to exceed the preceding incumbent's payoff upper endpoint. Increasing a lower endpoint by itself is insufficient. Rejected proposals would restore both the incumbent actor and optimizer state, and all proposals remain recorded. Every one of the fifteen noninitial proposals passes this strict value test.

Proposition~\ref{main-prop:r19accept} gives the payoff-ordering and observable regret-contraction implications. Theorem~\ref{main-thm:r19transfer} then certifies a positive initial-to-final true payoff gain uniformly on $K$, using common shocks and a difference-sensitive analytic modulus rather than point samples. The uniform gain lower bound is at least $\NineteenGain$. Combining gain and final regret gives true-regret reduction of at least 58 percent on $K$ for each seed. At the reference state the direct 75-percent test also passes. None of these statements can be produced merely by replacing a critic while holding the actor fixed.

We also perform the requested historical attribution audit for all ten accessible feedback seeds. In every row, the initial actor paired with the final critic reproduces the final total bound, and the final policy-side negative component is zero. The complete 40-object factorial appears in the main table. Its adverse interpretation is preserved; the new value-sensitive experiment does not rewrite that historical result.

\subsection{R18-F3 --- A policy-learning implication for the executed mechanism}
\emph{Response.} The new mechanism uses separate policy evaluation and optimal comparison. If $L_{\pi'}-U_\pi\geq\kappa(\overline V-L_\pi)$, then the actual regret contracts by at least $1-\kappa$. This statement follows from ordered true payoffs, not from an optimized envelope being numerically smaller. It is linked to the implemented actor gate and instantiated on all retained trajectories.

The explicit audit verifies $\kappa=.02$ for all fifteen checkpoint transitions and $\kappa=.75$ for each reference-state initial-to-final pair. On $K$, the alternative gain/regret criterion proves a final-to-initial true-regret ratio below $.42$. These are a posteriori learning statements for the executed policy sequence. We do not turn them into a claim that Adam converges for every initialization, that every proposed actor passes, or that the rate applies at unexecuted work levels.

The earlier complete-cover surrogate theorem is retained as a theorem about witness/certificate-aware proposal generation. Its role is not confused with the new policy-value acceptance theorem. The finite-MDP statements remain auxiliary, not an unexecuted continuous-time policy-iteration implementation.

\subsection{R18-F4 --- Continuation regime and the fixed-witness floor}
\emph{Response.} The new policies supply constructive continuation targets from their own independently evaluated original stopped payoffs. Freeze the generated controls, remove the initial-wealth transfer used on $K$, and let wealth approach 2 from the interior. Wealth is deterministic, strictly inward, and above the lower boundary until the preference stop or terminal date. The policy-value limit therefore exists. It is computed with the same original-utility quadrature and uniform stopping allowance, under separately checked no-transfer feasibility conditions.

Proposition~\ref{main-prop:r19trace} gives an excess of at least $\NineteenTrace$ above settlement, uniformly for initial preference in $[1.98,2.02]$ and for all five policies. This is a policy-generated continuation-value target in the correct regime, not a small arbitrary critic bias. It is explicitly an interior limit, not the value of a process initialized exactly on the absorbing boundary.

The target gives a constructive necessary trace constraint for an optimal upper witness. It is not itself a globally verified optimal upper witness. We test the tempting scalar-bias repair separately: increasing only each old critic's last bias enough to clear its pointwise necessary trace condition raises complete bounds to about 22.7--23.4. Those five adverse intervention outcomes are retained. The experiment shows why an interior residual solve must accompany trace correction. Thus R19 supplies a meaningful continuation target and a proof of its scope, but does not claim that bias lifting or a point-trace check alone solves the full-domain witness problem.

\subsection{R18-F5 --- Accessibility and finite-budget advantage}
\emph{Response.} The matched accessibility/all-face experiment remains unchanged. The main text states that width 16 and width 32 favor opposite variants. The structural result removes an inconsistent approximation-space constraint; it is not advertised as a uniform finite-budget advantage. The new abstract bases its numerical policy-learning claim on independently ordered payoffs, not on this architecture comparison. All prior boundary and accessibility theorems are retained.

\subsection{R18-F6 --- Matched certified accuracy and classical comparison}
\emph{Response.} We generate fresh SLSQP time-control transcriptions at 16, 32, and 64 slabs and evaluate them with the same independent stopped-payoff procedure and the same certified optimal upper comparator. All reach the reference-state $.01$ threshold. The neural and classical first-attainment frontier includes generation, compilation where applicable, and independent policy evaluation; preceding neural checkpoint checks are charged. Shared upper-comparator generation and verification costs are explicitly identified as historical costs rather than treated as zero or newly measured. Process-wide peak memory is retained.

The classical time transcription is faster in this low-dimensional control class. We do not conceal that outcome or claim a neural speed advantage. A further coupled-payoff comparison bounds the fresh 16-slab classical policy's advantage over each final neural policy by $\NineteenClassicalLoss$ uniformly on $K$, without assuming that the classical policy is optimal.

This supplies a sharp matched reference-state frontier and a direct uniform policy comparison. It does not establish a sharp full-domain Markov-chain or semi-Lagrangian frontier. Those six feedback candidates retain their original identities, actions, and complete certificates; all signed components are now reported together. The remaining global comparator task is not declared closed by a different numerical object.

\subsection{R18-F7 --- Identifying the learned high-dimensional contribution}
\emph{Response.} The new learned-gradient theorem depends explicitly on the initialized plan's gradient norm. It bounds corrected loss by $G_\theta(x)^2(145/209)^{2j}/8$. Verified initial-gradient maxima therefore produce learned-quality-dependent work budgets on each finite query set. A separate covering extension gives the precise additional Lipschitz and radius terms needed for a state-uniform conclusion. Thirty-two queries are not described as a cover of a 128-dimensional cube.

The new predeclared study compares all three stored neural initializers, zero-start gradient descent, and zero-start acceleration on 32 new queries per dimension at four common target tolerances. Every delivered plan receives its own directed certificate. It counts actual corrections to target, not the errors after an arbitrary common correction count. In 128 dimensions at tolerance $.01$, neural initialization uses five mean corrections versus six for zero-start descent and approximately 5.97 for acceleration. Each neural seed uses fewer gradients than acceleration on 31 of 32 queries. At stricter tolerances the classical accelerated method catches up or wins; those outcomes remain in the tables.

The contribution is a certified finite-query accuracy/work benefit in specified regimes, not a general neural necessity claim. All-seed break-even accounting is supplied with inference, correction, checking, and recorded training costs. Transferring historical training times is disclosed; no matched-hardware or end-to-end 128-dimensional speedup is asserted.

\subsection{R18-F8 --- Neural methodological identity}
\emph{Response.} R19 now supplies newly trained neural policies whose actual payoff improves and whose reference-state error is certified below $.01$ in the original economy. Their parameters are optimized only against the model payoff, not fitted to an inherited time-control library. The inherited dual is an independent comparison used after proposal generation. A common non-neural evaluator is not relabeled as a neural actor, just as using a conventional error estimator does not change the provenance of a numerical approximation.

The actor's time-control structure and state-transfer rule are stated prominently, so this new result does not silently claim general full-state feedback approximation. The paper keeps its title and neural-solver objective, while reporting separately what is established for the new policy block, the old feedback block, and the high-dimensional correction specialization. The strong non-neural historical library remains named as such and is not used as a fictitious new neural run.

\subsection{R18-F9 --- Difficulty and generality of the nonlinear application}
\emph{Response.} The nonlinear inventory application is retained in full, including its deterministic 12-date horizon, unconstrained controls, sparse derivatives, and global strong convexity. Its favorable classical structure is neither hidden nor rhetorically equated with a general stochastic curse-of-dimensionality benchmark. The new study makes the learned contribution testable within that application: target-based gradient work, learned initial-gradient bounds, and explicit amortization accounting.

The results establish a regime-dependent learned warm-start benefit, including a finite-query gradient-work gain at 128 dimensions and tolerance $.01$. They do not establish broad neural dominance on difficult stochastic dynamic models. We have not substituted an unexecuted harder model for evidence, and have not removed the existing model or its classical comparisons. The original stopped stochastic economy remains the primary policy-value setting, with its full-domain objective intact.

\subsection{R18-F10 --- Policy quality rather than descriptive controls}
\emph{Response.} The new evidence directly concerns original economic payoffs. It includes disjoint policy-value intervals at every checkpoint, a positive true-payoff gain throughout $K$, a true-regret contraction bound, a rectangular optimality bound, and a direct neural/classical payoff comparison on that same rectangle. These quantities address the requested nontrivial subset of the state domain and do not infer optimal controls from their ranges. Historical control ranges are preserved as feasibility/descriptive records, not reinterpreted as accuracy certificates.

\subsection{R18-F11 --- Economic meaning for the neural result}
\emph{Response.} The new neural policies have their own budget-financed sufficient compensation. Increase initial wealth by $.0092$ and add $.0092/Q_r$ to every consumption slab while holding preference and portfolio controls fixed. Directed checks establish the original consumption cap over $K$, identical terminal wealth, inward feasible wealth paths, and a common preference-exit time. Proposition~\ref{main-prop:r19wealth} gives a payoff gain larger than $.0119058$, exceeding every new rectangular neural regret bound.

The wealth increment is 0.736 percent of the reference wealth 1.25. It is sufficient compensation with an explicit feasible policy adjustment, not the exact minimum compensating variation of an unchanged policy. The earlier 0.624-percent non-neural compensation remains preserved and separately labeled. Neither number is used to normalize the old 7-unit full-domain feedback certificate.

\section{Technical comments T1--T10}
\subsection{R18-T1 --- Main-text signed decomposition}
The main factorial table now includes the positive and negative components alongside all four actor--critic combinations. It states directly that every old final feedback certificate has zero policy-side negative component and that its full remaining bound is the upper-comparison component. The new policy-value table is adjacent to the discussion of what constitutes actual policy improvement.

\subsection{R18-T2 --- All-seed interchange}
All ten accessible seeds have the complete initial/initial, final/initial, initial/final, and final/final comparison. The twenty cross-combinations are independently recomputed; the twenty unchanged diagonal records are identified as inherited checks. Their exact frozen networks and cell records remain available. Cross-seed critic swaps are not substituted for this requested within-seed factorial and are not claimed executed.

\subsection{R18-T3 --- Trace trajectories}
All thirty accessible checkpoints now have verified upper-face trace excess, the implied fixed-witness floor, and the positive comparison component in two complete supplement tables. The separate five bias-only interventions retain both adverse complete certificates and their achieved point-trace levels. The constructive new-policy continuation evaluations have their own files and are not mixed with old-critic trajectories.

\subsection{R18-T4 --- Acceptance terminology}
The complete-cover algorithm is called certificate-incumbent replacement, not policy improvement merely because its bound descends. Actual policy improvement in R19 refers only to independently ordered policy payoffs under the new gate or the coupled-payoff theorem. Its observable contraction proposition states exactly the additional verified inequality required.

\subsection{R18-T5 --- Scaling of complete-cover work}
The main text now expresses direct dense-jet work in the number of boxes, state dimension, network width/depth, and action-maximization cost. It reports the derivative-component count, intermediate-memory order, the effect of chunked verification, the tensor-cover product, and the unresolved costs of interval dependency/subdivision. This accounting does not claim a dimension-free complete-cover method. The sparse inventory correction has a separate $O(Hd)$ cost statement.

\subsection{R18-T6 --- Initialization-agnostic guarantee}
The high-dimensional discussion begins by stating that the previous cube-uniform theorem holds for an untrained bounded initializer and does not use learned quality. The new gradient-sensitive theorem is a separate statement with a separate numerical instantiation and explicit cover requirements. Their guarantees are not conflated.

\subsection{R18-T7 --- Tolerance and amortization tables}
Every dimension, neural seed, and target $10^{-2},10^{-3},10^{-4},10^{-6}$ has target-attainment counts, paired zero/accelerated counts, inference/correction time, final checking time, and break-even accounting. No finite break-even is returned for nonpositive measured savings. Solver-only and check-inclusive accounting are distinguished, and historical training-time transfer is disclosed. All 1,920 final target-check records and exact stored plans are retained. Work certificates concern the delivered plans; measured wall times are not mathematical speedup certificates.

\subsection{R18-T8 --- Classical finest-grid deterioration}
The new signed table covers all six original feedback baselines. At the finest grids the policy-side term becomes positive while the witness-side term changes comparatively little. The main text states both the arithmetic diagnosis and its limit: identifying a residual component is not equivalent to proving lower true policy quality or isolating interpolation from evaluation-witness mismatch. No unfavorable row is removed.

\subsection{R18-T9 --- Arithmetic versus mathematics}
The distinction is preserved throughout. MPFR-directed arithmetic provides an independent arithmetic path for the checked formulas, whereas model derivations, quadrature mathematics, stopping inequalities, and comparison logic remain shared mathematical premises. Primitive tests, file hashes, and build checks are not described as an independent formal proof of the entire theorem stack. New analytic arguments are written out in the main appendix for substantive scrutiny.

\subsection{R18-T10 --- Prospective execution}
The R19 protocol predates the new neural training and finite-query study. It fixes seeds, work levels, architecture, proposal quadrature, acceptance semantics, target tolerances, correction cap, comparators, and retention of failures. The final reproduction enforces the policy-value gate during training rather than merely attaching it to an old trajectory afterward. Deterministic state-extension, trace, and compensation proofs are identified as subsequent analytic consequences. Historical R17 swaps remain an explicitly retrospective audit, not new randomized R19 training.

\section{Status of the requested R19 gates}
Gate A now has policy-sensitive progress in the unchanged economy: actual accepted payoff improvements at the reference state and a proved uniform gain and true-regret reduction on $K$. Gate B now has constructive feasible-policy continuation targets in the correct regime; a globally sharp optimal upper witness for the full feedback cylinder is still not supplied. Gate C has new sub-$.01$ reference-state neural accuracy and a $\NineteenK$ rectangular bound, but the entire-cylinder $.01$ target remains unmet. Gate D has a sharp matched reference-state time-control frontier and direct uniform policy comparisons; a comparably sharp global Markov-chain/semi-Lagrangian frontier remains unestablished. Gate E has learned-quality-dependent work bounds and finite-query gradient-work benefits at stated tolerances, not a universal high-dimensional speedup. Gate F is supported by new neural policy-learning evidence without renaming non-neural successes or deleting the former developments.

The revision is submitted for substantive reassessment of these new objects and proofs. It does not equate successful arithmetic tests or PDF compilation with closure of all scientific gates. The review entry point identifies the exact manuscripts, complete results, reproduction commands, and preservation manifest. Original review branches and prior manuscript files are left unchanged.
